\chapter{Review of calculus}
\label{ch:Lec-1-calculus}

The following is a brief review of some concepts from 
calculus that we will use throughout the course. 
For a more detailed treatment see the classic texts 
\cite{apostol1991calculus-1,apostol1991calculus-2, babyrudin}.


\section{Partial derivatives, gradients, and other useful operators}

Throughout this section we consider a scalar function $u: \Omega \to \R$
or a vector-valued function $\fb: \Omega \to \R^m$, where $\Omega
\subseteq \R^d$ is an open set and $\xb = (x_1, x_2, \ldots, x_d) \in
\Omega$. We write $u \in C^k(\Omega)$ to mean that $u$ is $k$-times
continuously differentiable on $\Omega$.

\begin{definition}[Partial derivative]
    The partial derivative of $u$ with respect to $x_i$ at a point
    $\xb \in \Omega$ is defined as
    \begin{equation*}
        \frac{\partial u}{\partial x_i}(\xb) :=
        \lim_{h \to 0} \frac{u(\xb + h \mbf{e}_i) - u(\xb)}{h},
    \end{equation*}
    provided the limit exists, where $\mbf{e}_i \in \R^d$ is the $i$-th
    standard basis vector.
\end{definition}
In words, $\partial u/\partial x_i$ measures the rate of change of $u$
as we move in the $x_i$ direction while holding all other coordinates
fixed. More generally, the directional derivative of $u$ at $\xb$ in the
direction $\vb \in \R^d$ is defined as
\begin{equation*}
    D_\vb u(\xb) := \lim_{h \to 0} \frac{u(\xb + h\vb) - u(\xb)}{h},
\end{equation*}
which reduces to the partial derivative $\partial u/\partial x_i$ when
$\vb = \mbf{e}_i$.

\begin{definition}[Gradient]
    Suppose $u \in C^1(\Omega)$. The gradient of $u$ at $\xb \in \Omega$
    is the vector of first partial derivatives
    \begin{equation*}
        \nabla u(\xb) :=
        \begin{bmatrix}
            \dfrac{\partial u}{\partial x_1}(\xb) \\[1ex]
            \vdots \\[0.5ex]
            \dfrac{\partial u}{\partial x_d}(\xb)
        \end{bmatrix} \in \R^d.
    \end{equation*}
\end{definition}
The gradient encodes all directional derivatives of $u$ at once via the
identity $D_\vb u(\xb) = \langle \nabla u(\xb), \vb \rangle$, and it
points in the direction of steepest increase of $u$ at $\xb$; we will
use this fact later in the course when we discuss optimization-based
methods.

\begin{lemma}[Equality of mixed partials]
    \label{lem:mixed-partials}
    If $u \in C^2(\Omega)$, then
    \begin{equation*}
        \frac{\partial^2 u}{\partial x_i \partial x_j}(\xb) =
        \frac{\partial^2 u}{\partial x_j \partial x_i}(\xb),
        \qquad \text{for all } \xb \in \Omega, \; i,j = 1,2,\ldots,d.
    \end{equation*}
\end{lemma}
In other words, as long as $u$ is twice continuously differentiable, the
order in which we take mixed partial derivatives does not matter.

\begin{definition}[Jacobian]\label{def:jacobian}
    Suppose $\fb = (f_1, f_2, \ldots, f_m): \Omega \to \R^m$ with each
    $f_i \in C^1(\Omega)$. The Jacobian of $\fb$ at $\xb \in \Omega$ is
    the matrix of first partial derivatives
    \begin{equation*}
        D\fb(\xb) :=
        \begin{bmatrix}
            \nabla f_1(\xb)^\top \\
            \vdots \\
            \nabla f_m(\xb)^\top
        \end{bmatrix}
        =
        \left(
            \frac{\partial f_i}{\partial x_j}(\xb)
        \right)_{i=1,\ldots,m}^{j=1,\ldots,d}
        \in \R^{m \times d}.
    \end{equation*}
\end{definition}
When $m=1$ the Jacobian is just the (transposed) gradient, i.e.,
$Du(\xb) = \nabla u(\xb)^\top$. The Jacobian is the natural
generalization of the derivative to vector-valued functions of several
variables: it is the matrix representing the best linear approximation
of $\fb$ near $\xb$, i.e., $\fb(\xb + \mbf{h}) \approx \fb(\xb) +
D\fb(\xb) \mbf{h}$ for $\mbf{h} \in \R^d$ small.

\begin{definition}[Hessian]
    Suppose $u \in C^2(\Omega)$. The Hessian of $u$ at $\xb \in \Omega$
    is the matrix of second partial derivatives
    \begin{equation*}
        D^2 u(\xb) :=
        \left(
            \frac{\partial^2 u}{\partial x_i \partial x_j}(\xb)
        \right)_{i,j=1}^d \in \R^{d \times d}.
    \end{equation*}
\end{definition}
Note that $D^2 u(\xb) = D(\nabla u)(\xb)$, i.e., the Hessian is the
Jacobian of the gradient. By the equality of mixed partials, $D^2
u(\xb)$ is a symmetric matrix whenever $u \in C^2(\Omega)$.

Finally, we introduce two operators acting on vector fields and scalar
fields, respectively, that will reappear frequently once we start
discussing PDEs.

\begin{definition}[Divergence]
    Suppose $\fb = (f_1, \ldots, f_d): \Omega \to \R^d$ with each $f_i
    \in C^1(\Omega)$. The divergence of $\fb$ at $\xb \in \Omega$ is the
    scalar
    \begin{equation*}
        \nabla \cdot \fb(\xb) :=
        \sum_{i=1}^d \frac{\partial f_i}{\partial x_i}(\xb) =
        \mathrm{tr}\big(D\fb(\xb)\big).
    \end{equation*}
\end{definition}

\begin{definition}[Laplacian]
    Suppose $u \in C^2(\Omega)$. The Laplacian of $u$ at $\xb \in
    \Omega$ is the scalar
    \begin{equation*}
        \Delta u(\xb) := \nabla \cdot \nabla u(\xb) =
        \sum_{i=1}^d \frac{\partial^2 u}{\partial x_i^2}(\xb).
    \end{equation*}
\end{definition}
The Laplacian already made an appearance in \Cref{ex:heat_equation}: in
one spatial dimension it is simply $\Delta u = \partial^2 u/\partial
x^2$, and the heat equation reads $\partial u/\partial t = \alpha \Delta
u$. We will see many more PDEs written in terms of the Laplacian and
other differential operators introduced above later in the course.


\section{Taylor's theorem}

Taylor's theorem is arguably the single most important tool from
calculus for this course: nearly every numerical method we study is
derived, and its accuracy analyzed, by Taylor expanding the exact
solution and comparing it to the numerical approximation.

\begin{theorem}[Taylor's theorem, one variable]
    \label{thm:taylor-1d}
    Suppose $u \in C^{p+1}([a,b])$ for some integer $p \ge 0$ and let
    $x, x+h \in [a,b]$. Then
    \begin{equation*}
        u(x+h) = \sum_{k=0}^p \frac{u^{(k)}(x)}{k!} h^k +
        \frac{u^{(p+1)}(\xi)}{(p+1)!} h^{p+1},
    \end{equation*}
    for some $\xi$ between $x$ and $x+h$.
\end{theorem}
The polynomial $\sum_{k=0}^p u^{(k)}(x) h^k / k!$ is called the $p$-th
order Taylor polynomial of $u$ about $x$, and the last term is the
remainder; the theorem says that the remainder is exactly the "next"
term in the series, evaluated at some unknown intermediate point $\xi$
rather than at $x$ itself. In numerical analysis it is often more
convenient to bound the size of the remainder rather than track it
exactly, which motivates the following standard notation.

\begin{definition}[Big-$\mO$ notation]
    We say $f(h) = \mO(g(h))$ as $h \to 0$ if there exist constants $C >
    0$ and $\delta > 0$ such that $|f(h)| \le C |g(h)|$ for all $|h| <
    \delta$.
\end{definition}
For example, if $u \in C^{p+1}$ near $x$ with $|u^{(p+1)}|$ bounded,
\Cref{thm:taylor-1d} implies
\begin{equation*}
    u(x+h) = \sum_{k=0}^p \frac{u^{(k)}(x)}{k!} h^k + \mO(h^{p+1}),
\end{equation*}
which is the form in which we will use Taylor's theorem most often.

Taylor's theorem also extends to functions of several variables; we
state the second order version here since it is the one we will use
most.

\begin{theorem}[Taylor's theorem, several variables]
    Suppose $u \in C^3(\Omega)$ for $\Omega \subseteq \R^d$ open and let
    $\xb, \xb + \mbf{h} \in \Omega$ with the line segment connecting them
    contained in $\Omega$. Then
    \begin{equation*}
        u(\xb + \mbf{h}) = u(\xb) + \langle \nabla u(\xb), \mbf{h} \rangle +
        \frac12 \langle D^2 u(\xb) \mbf{h}, \mbf{h} \rangle + \mO(\|\mbf{h}\|_2^3),
    \end{equation*}
    as $\mbf{h} \to 0$, where $\nabla u(\xb)$ and $D^2 u(\xb)$ are the
    gradient and Hessian of $u$ at $\xb$, respectively.
\end{theorem}
This is simply the vector generalization of the familiar quadratic
Taylor approximation from single-variable calculus, with the gradient
playing the role of $u'(x)$ and the Hessian playing the role of
$u''(x)$.


\section{Integration by parts and Green's identities}

Just as Taylor's theorem is our main tool for analyzing numerical
differentiation, integration by parts is our main tool for making sense
of weak or variational formulations of PDEs, which underlie finite
element methods discussed later in the course.

\begin{theorem}[Integration by parts]
    Suppose $u, v \in C^1([a,b])$. Then
    \begin{equation*}
        \int_a^b u'(x) v(x) \, dx = u(b)v(b) - u(a)v(a) -
        \int_a^b u(x) v'(x) \, dx.
    \end{equation*}
\end{theorem}
In higher dimensions the analogous tool is the divergence theorem,
which relates an integral over a domain to an integral over its
boundary.

\begin{theorem}[Divergence theorem]
    Let $\Omega \subset \R^d$ be a bounded domain with piecewise smooth
    boundary $\partial \Omega$ and outward unit normal $\mbf{n}$, and
    suppose $\fb \in C^1(\overline\Omega)^d$. Then
    \begin{equation*}
        \int_\Omega \nabla \cdot \fb(\xb) \, d\xb =
        \int_{\partial \Omega} \langle \fb(\xb), \mbf{n}(\xb) \rangle \, dS(\xb).
    \end{equation*}
\end{theorem}
Applying the divergence theorem to the vector field $\fb = v \nabla u$
for scalar functions $u \in C^2(\overline\Omega)$ and $v \in
C^1(\overline\Omega)$, and using the product rule $\nabla \cdot (v
\nabla u) = v \Delta u + \langle \nabla u, \nabla v \rangle$, yields the
following pair of identities.

\begin{corollary}[Green's first identity]
    \label{cor:greens-first-identity}
    Under the assumptions above,
    \begin{equation*}
        \int_\Omega \langle \nabla u(\xb), \nabla v(\xb) \rangle \, d\xb
        = - \int_\Omega v(\xb) \Delta u(\xb) \, d\xb +
        \int_{\partial \Omega} v(\xb) \frac{\partial u}{\partial n}(\xb)
        \, dS(\xb),
    \end{equation*}
    where $\partial u/\partial n := \langle \nabla u, \mbf{n} \rangle$
    is the normal derivative of $u$ on $\partial \Omega$.
\end{corollary}

\begin{corollary}[Green's second identity]
    If in addition $v \in C^2(\overline\Omega)$, then
    \begin{equation*}
        \int_\Omega \big( u(\xb) \Delta v(\xb) - v(\xb) \Delta u(\xb)
        \big) \, d\xb =
        \int_{\partial \Omega} \left(
            u(\xb) \frac{\partial v}{\partial n}(\xb) -
            v(\xb) \frac{\partial u}{\partial n}(\xb)
        \right) dS(\xb).
    \end{equation*}
\end{corollary}
Green's first identity is the key step in deriving the weak formulation
of PDEs such as Poisson's equation from \Cref{ex:poisson-bvp}.


\section{Differentiating functions of vectors and matrices}

We close this review by collecting a few differentiation rules for
functions of vectors and matrices that will repeatedly come in handy,
for instance when linearizing systems of nonlinear ODEs or setting up
implicit time-stepping schemes.

\begin{lemma}[Chain rule]
    Suppose $\fb: \Omega \to \R^m$ and $\gb: U \to \R^p$ with $\fb(\Omega)
    \subseteq U \subseteq \R^m$, $\fb \in C^1(\Omega)$, and $\gb \in
    C^1(U)$. Then $\gb \circ \fb: \Omega \to \R^p$ is differentiable and
    \begin{equation*}
        D(\gb \circ \fb)(\xb) = D\gb(\fb(\xb)) \, D\fb(\xb),
        \qquad \xb \in \Omega,
    \end{equation*}
    where the right-hand side is a product of a $p \times m$ matrix with
    an $m \times d$ matrix.
\end{lemma}

The following gradient and Hessian formulas for linear and quadratic
functions of a vector $\xb \in \R^d$ are especially useful and worth
memorizing; see also \cite{matrixcookbook} for a much more extensive
list of matrix calculus identities.

\begin{itemize}
    \item For a fixed vector $\mbf{a} \in \R^d$ and $u(\xb) = \langle \mbf{a},
    \xb \rangle$, we have $\nabla u(\xb) = \mbf{a}$ and $D^2 u(\xb) = 0$.
    \item For a fixed matrix $A \in \R^{d \times d}$ and $u(\xb) =
    \langle A \xb, \xb \rangle$, we have $\nabla u(\xb) = (A + A^\top)
    \xb$ and $D^2 u(\xb) = A + A^\top$; if $A$ is symmetric this
    simplifies to $\nabla u(\xb) = 2A\xb$ and $D^2 u(\xb) = 2A$.
    \item In particular, for $u(\xb) = \|\xb\|_2^2 = \langle \xb, \xb
    \rangle$, we have $\nabla u(\xb) = 2\xb$ and $D^2 u(\xb) = 2I_d$.
\end{itemize}

Finally, we will also need to differentiate matrix-valued functions of
a single scalar variable, such as a time-dependent matrix $A(t) \in
\R^{n \times n}$. This is done entrywise, i.e., $A'(t) := (a_{ij}'(t))$,
and the usual product rule carries over:
\begin{equation*}
    \frac{d}{dt} \big( A(t) B(t) \big) = A'(t) B(t) + A(t) B'(t).
\end{equation*}
A particularly important example for us is the matrix exponential of a
constant matrix $A \in \R^{n \times n}$ introduced in
\Cref{ch:Lec0-linalg}, which satisfies
\begin{equation*}
    \frac{d}{dt} e^{tA} = A e^{tA} = e^{tA} A, \qquad t \in \R.
\end{equation*}
This identity is what makes the matrix exponential the natural building
block for solutions of linear systems of ODEs with constant
coefficients.

\bibliographystyle{plainnat}
\bibliography{references}